\documentclass[11pt]{article}
\usepackage{setspace}
\usepackage{geometry}
\geometry{margin=1in}
\setstretch{1.15}

\begin{document}

% --- Applicant Information ---
\noindent
\textbf{Zhongli Wang} \\
Phone: +49 160 9541 6015 \\
Email: nkwangzhongli@gmail.com \\
Website: https://sites.google.com/view/zhongli-wang \\
Bielefeld, Germany \\

\vspace{0.5cm}

% --- Recipient Information ---
\noindent
\textbf{Recruitment Committee} \\
Department of Economics \\
Copenhagen Business School \\
Copenhagen, Denmark \\

\vspace{0.5cm}

% --- Subject Line ---
\noindent
\textbf{Application for the Postdoc Position in Behavioral and Applied Microeconomics} \\

\vspace{0.3cm}

\noindent Dear Members of the Recruitment Committee,\\[0.2cm]

I am writing to apply for the Postdoc in Behavioral and Applied Microeconomics at Copenhagen Business School. I am completing my Ph.D.\ in Economics at Bielefeld University and Università Cattolica del Sacro Cuore through the Marie Skłodowska–Curie EPOC programme. My research combines applied microeconomics, behavioral modeling, and large-scale computational methods to study belief formation, decision-making under uncertainty, and the behavioral consequences of new technologies. These interests align strongly with the research agendas led by Associate Professor Luigi Butera and Professor Dmitry Taubinsky.

I began my Ph.D.\ in 2021, during a period when the COVID-19 pandemic accelerated digitalization and exposed how technological shocks can rapidly alter markets and individual behavior. This context motivated the central chapter of my dissertation, which examines how firms in ICT markets adjust capacity, pricing, and quality as technologies evolve and user behavior adapts. Using firm-level data, I estimate key structural parameters and evaluate counterfactual technology-adoption scenarios through simulation-based econometrics. My dissertation, \textit{“Strategic Behavior and Economic Dynamics in ICT Markets: An Agent-Based Framework,”} develops computational models in which firms and heterogeneous users respond to information frictions, congestion, and shifting technological conditions. To analyze these dynamics, I combine Monte Carlo simulation, neural-network surrogate models, and structural estimation to study how innovations diffuse and how strategic conduct changes under alternative technological and regulatory environments.

A unifying theme across my doctoral work is understanding how individuals form expectations, update beliefs, and adapt their behavior in complex or uncertain environments. I develop empirical, experimental, and simulation-based models to study heterogeneous behavioral responses, misperceptions, and learning under imperfect information. These tools naturally complement the CBS research team’s agenda on subjective well-being, financial expectations, and the role of memory and beliefs in economic decisions. My background in modeling heterogeneity, belief formation, and behavioral dynamics equips me to contribute directly to large-scale survey--administrative data projects investigating biases, forecasting errors, and emotionally driven economic choices.

In addition to applied microeconometrics, I bring experience working with high-dimensional data, behavioral simulations, and AI-assisted empirical analysis. One coauthored project, \emph{Polytope Fraud Theory}, uses machine learning to detect behavioral anomalies in large financial datasets, illustrating my ability to merge predictive modeling with economic interpretation. A second project develops an agentic workflow integrating large-language-model agents into empirical pipelines, offering new ways to analyze beliefs, linguistic patterns, and cognitive biases at scale. These skills would complement the team’s plan to merge large survey and choice experiments with rich administrative data and to study heterogeneity in expectations, emotional shocks, and behavioral errors.

I am particularly excited about the opportunity to contribute to the three research agendas outlined in the call:
\begin{itemize}
    \item studying determinants of subjective well-being using large-scale survey–administrative data linkages;
    \item analyzing biases and misperceptions in individuals’ financial expectations and their behavioral implications;
    \item investigating how short-lived emotional shocks affect financial and health decisions and how memory shapes these effects.
\end{itemize}
My prior work on belief formation, dynamic adjustment, learning under uncertainty, and heterogeneous behavioral types provides a strong methodological base for these projects. I would be enthusiastic to help build new datasets, design survey and experimental instruments, develop empirical strategies, and coauthor publications with the CBS–Berkeley research team.

Beyond technical skills, I greatly value collaborative research environments. The prospect of working closely with Professors Butera and Taubinsky and an interdisciplinary team of students and RAs is a major motivation for me. I would welcome the opportunity to contribute actively to data construction, experiment design, manuscript preparation, and the broader intellectual life of the department.

Thank you very much for considering my application. I would be delighted to contribute to CBS’s behavioral and applied microeconomics group and to collaborate on high-impact research that advances our understanding of beliefs, biases, and economic decision-making. I would welcome the opportunity to discuss my candidacy further.

\vspace{6mm}

\noindent Sincerely,\\[15mm]
\noindent Zhongli Wang

\end{document}
